% Copulas & Tail Dependence — research-note PDF.
%
% PROSE AND CODE ARE VERBATIM from the website article:
%   site/src/content/articles/copulas-tail-dependence.tsx
% The article interpolates its computed results into the prose; every
% interpolated value here is the one the site renders, read from
% data/copulas-tail-dependence.ts. Code cells are inline in the article.
%
% CHARTS generated by figures/make_copulas_tail_dependence.py from that module.
\documentclass{pyportfolios-note}

\noteshorttitle{Copulas \& Tail Dependence}

\begin{document}
\thispagestyle{notecover}

\notemasthead{Quantitative Insights}{Tutorial}{Q3 2026}{}

\notetitle{Copulas \&\\Tail Dependence}

\notedeck{Correlation is a photograph taken in calm weather. The corner of the
joint distribution is where the money is lost.}

\notelede{A correlation matrix is a photograph of dependence taken in calm
weather. The question risk managers actually get paid to answer — \textbf{if the
S\&P has one of its worst weeks, what happens to London and Tokyo?} — lives in
the corner of the joint distribution, where a single number cannot see. Copulas
separate what each market does on its own from how they move together, and once
you make that split on 25 years of real index data, the verdict is blunt: the
Gaussian copula says joint crashes are vanishingly rare; the data says they
cluster; and a Student-t copula with 3.5 degrees of freedom repairs most of the
damage.}

\notecard
  {Risk Management}
  {SciPy · NumPy · Pandas}
  {\^{}GSPC, \^{}FTSE, \^{}N225}
  {Jan 2000 – Dec 2024}
  {Stress-testing a multi-asset book against the joint tail, not the average}

\begin{multicols}{2}
\begin{notepipeline}
\item Weekly log returns for \^{}GSPC, \^{}FTSE, \^{}N225, 2000–2024
      (holiday-aligned)
\item Pearson vs Spearman vs Kendall — three answers to `how correlated?'
\item Rank-transform to pseudo-observations on the unit square
\item Fit Gaussian and Student-t copulas (df by MLE over a 2–30 grid)
\item Measure empirical tail dependence \& conditional crash probabilities
\end{notepipeline}
\end{multicols}

% =========================================== 1 Dependence is not correlation =
\noteneedroom{190bp}
\notesection{1.\notenumsep Dependence is not correlation}

The sample is 1,302 aligned weekly returns for the S\&P 500, FTSE 100 and Nikkei
225, 2000-01-03 to 2024-12-31 — dot-com bust, 2008, the euro crisis, COVID and
the 2022 rates shock all included. Before any statistics, two traps come free
with global indices, and both are worth learning to spot. First, each exchange
keeps its own holiday calendar, so returns are computed jointly and any week
missing a market is dropped. Second — subtler — the markets do not even trade at
the same time. Tokyo closes roughly ten hours before London and some fifteen
before New York, so a same-calendar-day Nikkei close \textbf{leads} the US close,
and daily cross-market correlations are structurally understated: the two prices
are answering questions asked at different moments. Resampling to weekly
(Friday-to-Friday) returns absorbs most of that offset — watch the SPX–Nikkei
Pearson correlation jump from \texttt{0.15} on daily data to \texttt{0.58} on
weekly.

Even then, ``the correlation'' is three different numbers. Pearson measures
linear co-movement and is hostage to outliers; Spearman and Kendall are
rank-based, so they see only the ordering — exactly the part a copula models.
Kendall's τ has the cleanest interpretation: the probability that two randomly
chosen weeks agree in direction, minus the probability they disagree.

\notefigure{figures/copulas-tail-dependence/fig1-kendall.png}
  {Kendall's τ, weekly returns 2000–2024}
  {Rank correlation: SPX–FTSE 0.54, SPX–NKY 0.37, FTSE–NKY 0.35}
  {Source: Author calculations, from the article's computed results.}

\notecallout{Why practitioners care}{Every multi-asset stress scenario is
implicitly a statement about joint tails. If your scenario engine assumes a
dependence structure with no tail clustering, diversification will look better on
paper than it behaves in a crisis — precisely when you are counting on it.}

% ================================ 2 Sklar's theorem & pseudo-observations ====
\noteneedroom{190bp}
\notesection{2.\notenumsep Sklar's theorem \& pseudo-observations}

Here is the idea that makes the whole subject click. Sklar's theorem (1959) says
any joint distribution factors cleanly into two parts: the \textbf{marginals} —
what each market does on its own — and a \textbf{copula}, a joint distribution on
the unit square that carries all of the togetherness and none of the marginal
shape. Separate the dancers from the dance. To see the copula empirically,
replace each return by its normalised rank,
$u = \frac{\operatorname{rank}(x)}{n + 1}$ — the probability integral transform
done with the empirical CDF. Ranks are uniform by construction, so the marginals
are washed out entirely — any structure that survives in the picture below is
pure dependence.

\begin{notecode}
from scipy import stats

def pseudo_obs(x):                       # rank / probability integral transform
    return stats.rankdata(x) / (len(x) + 1)

U = pd.DataFrame({c: pseudo_obs(ret[c].values) for c in ret.columns})
\end{notecode}

\notefigure{figures/copulas-tail-dependence/fig2-pseudo.png}
  {Pseudo-observations, SPX vs FTSE — 400 of 1,302 weeks}
  {One point is one week; the dashed lines mark the 5\% quantile of each margin}
  {Source: Author calculations, from the article's computed results.}

If the two markets were independent this square would be filled uniformly.
Instead the mass drains toward the diagonal — and, critically, \textbf{piles up
in the corners}. The lower-left corner is the object of study for the rest of
this article: weeks in which both markets were simultaneously in their worst
tail.

% ===================== 3 The Gaussian copula — elegant, wrong in the corner ==
\noteneedroom{190bp}
\notesection{3.\notenumsep The Gaussian copula — elegant, and wrong in the corner}

\begin{multicols}{2}
The Gaussian copula is what you implicitly assume whenever you summarise joint
behaviour with a correlation matrix alone. Fitting it is one line: push the
pseudo-observations through the standard normal quantile function (``normal
scores'') and take their correlation — for SPX–FTSE that gives $\rho = 0.74$. It
is analytically convenient, it scales to any dimension, and it has one fatal
property: \textbf{zero tail dependence}. Ask it ``given one market is having a
1-in-q disaster, how often is the other one too?'' and as q deepens its answer
marches to zero — for any ρ < 1, no matter how high. In the limit, joint crashes
are not just rare. They are assumed away.

You can watch the assumption fail at finite depth. At the 10\% level the fitted
Gaussian copula implies a joint-crash ratio of 50.8\% for SPX–FTSE; at 5\% it has
slipped to 43.5\%; at 1\% it is down to 31.2\%, on its way to zero. The empirical
series goes the other way: 55.3\% at 10\%, 55.3\% at 5\%, 61.4\% at 1\%. The
deeper you look into the tail, the more the data diverges from the model — in the
direction that hurts.
\end{multicols}

\notecallout{The cautionary tale}{David Li's 2000 paper ``On Default Correlation:
A Copula Function Approach'' put the Gaussian copula at the centre of CDO
pricing, and for a few years a single correlation number priced trillions in
structured credit. The formula did exactly what it says here: it made
simultaneous defaults look almost impossible at depth. 2008 was, among other
things, the empirical tail-dependence estimate arriving.}

% ============================================= 4 The t copula & tail dep. ====
\noteneedroom{190bp}
\notesection{4.\notenumsep The t copula \& tail dependence}

The Student-t copula adds exactly one parameter — the degrees of freedom ν — and
that single knob buys tail dependence. We set ρ by Kendall's τ inversion,
$\rho = \sin\!\left(\tfrac{\pi\tau}{2}\right)$, which is exact for elliptical
copulas, and profile the exact copula log-likelihood over a ν grid from 2 to 30
on the SPX–FTSE pair:

\begin{notecode}
rho  = np.sin(np.pi * tau / 2)                  # tau inversion: 0.746
grid = np.arange(2.0, 30.5, 0.5)
ll   = np.array([t_copula_loglik(u, v, rho, df) for df in grid])
df_hat = grid[ll.argmax()]                      # -> 3.5

# lower-tail dependence, closed form (Demarta & McNeil 2005)
lam = 2 * stats.t.cdf(-np.sqrt((df_hat+1)*(1-rho)/(1+rho)), df=df_hat+1)
\end{notecode}

The likelihood picks $\nu = 3.5$ — heavy joint tails — and prefers the t copula
decisively: log-likelihood 577.9 against the Gaussian's 526.8 on the same
pseudo-observations, one extra parameter. Unlike the Gaussian, the t copula's
tail dependence does not vanish: it converges to the closed-form λ above, which
at ν = 3.5 gives 45.9\% for SPX–FTSE, 30.9\% for SPX–NKY and 29.2\% for
FTSE–NKY.

\notefigure{figures/copulas-tail-dependence/fig3-jointcrash.png}
  {Joint-crash ratio P(both < 5\% quantile) / 5\% — data vs models}
  {The Gaussian bar undershoots the observed ratio for every pair; the t-copula
   λ is a floor that does not decay}
  {Source: Author calculations, from the article's computed results.}

Read the empirical-vs-Gaussian gap per pair: the Gaussian copula, fitted to the
\textbf{same data} with the \textbf{same correlation}, undershoots the observed
5\% joint-crash ratio for every pair — and remember its bar keeps falling as q
shrinks while the empirical one rises. The middle bar is the t copula's limiting
λ: a floor that does not decay, sitting close to what the data shows.

% ============================== 5 Conditional crash probabilities ===========
\noteneedroom{190bp}
\notesection{5.\notenumsep Conditional crash probabilities}

The same mathematics, phrased the way a risk committee asks it: given that one
market has a worst-decile week, what is the probability the other one does too?
Under independence the answer would be 10\%. The data answers between 37.7\% and
55.4\%:

\notefigure{figures/copulas-tail-dependence/fig4-condcrash.png}
  {P(both in worst decile | one is) per pair — dashed line = the 10\% an independence assumption would give}
  {Every pair sits far above the independence line}
  {Source: Author calculations, from the article's computed results.}

\begin{notetable}{@{}p{110bp}p{190bp}>{\raggedleft\arraybackslash}p{90bp}>{\raggedleft\arraybackslash}p{100bp}@{}}
\thcell{Pair} & \thcell{P(both worst-decile | one is)} & \thcell{Gaussian @10\%} & \thcell{t copula λ (ν = 3.5)} \\
\theadrule
SPX–FTSE & 55.4\%  (72/130) & 50.8\% & 45.9\% \\ \rowrule
SPX–NKY & 42.3\%  (55/130) & 35.6\% & 30.9\% \\ \rowrule
FTSE–NKY & 37.7\%  (49/130) & 34.1\% & 29.2\% \\
\end{notetable}

A worst-decile S\&P week drags the FTSE into its own worst decile more than half
the time — 72 of 130 of the 130 conditioning weeks — and even the geographically
and temporally distant Nikkei follows 42.3\% of the time. Note the ordering in
every row: empirical ≥ t copula λ > Gaussian-at-depth. The single-parameter fix
gets you most, not all, of the way — the residual gap is the asymmetry between
crash and boom corners that symmetric elliptical copulas cannot express.

% ================================================= 6 The practitioner take ==
\noteneedroom{240bp}
\notesection{6.\notenumsep The practitioner take}

For multi-asset stress testing, the workflow this tutorial rehearses is the one
that survives contact with a crisis:

\begin{multicols}{2}
\begin{notebullets}
\item \textbf{Model marginals separately} — fat tails, volatility clustering,
      whatever each series needs — then choose the dependence structure as a
      deliberate act, not as a side effect of writing down a correlation matrix.
\item \textbf{Check the candidate copula against the data} — compare it to the
      empirical joint-crash ratios at several depths before trusting it.
\item \textbf{When in doubt between Gaussian and t, pay the one parameter} — the
      cheapest tail insurance in the toolbox: here it turned ``joint crashes
      become impossible'' into λ ≈ 46\% for the closest pair.
\item \textbf{Know where this road ends} — the t copula is elliptical, so its
      lower-tail dependence comes bundled with an identical upper tail: joint
      booms priced as generously as joint crashes. When that asymmetry matters,
      the next tools up are the asymmetric Archimedean families (Clayton glues
      lower tails only) and vine constructions, which assemble a
      high-dimensional copula from freely chosen pairs.
\end{notebullets}
\end{multicols}

\notecallout{Practitioner take}{Correlation tells you how markets dance; the
copula tells you how they die. Simulate scenarios from a fitted t copula
(ν = 3.5 here) rather than a Gaussian one and your stress engine will generate
the co-crashing weeks that actually populate the historical record — 130
worst-decile S\&P weeks in this sample, the majority of them shared with London.}

\noteneedroom{200bp}
\begin{notereferences}
\item Sklar, A. (1959). Fonctions de répartition à n dimensions et leurs marges.
      \textit{Publications de l'Institut de Statistique de l'Université de
      Paris}, 8, 229–231.
\item Embrechts, P., McNeil, A. \& Straumann, D. (2002). Correlation and
      dependence in risk management: properties and pitfalls. In \textit{Risk
      Management: Value at Risk and Beyond}, Cambridge University Press.
\item Demarta, S. \& McNeil, A. J. (2005). The t copula and related copulas.
      \textit{International Statistical Review}, 73(1), 111–129.
\item Aas, K., Czado, C., Frigessi, A. \& Bakken, H. (2009). Pair-copula
      constructions of multiple dependence. \textit{Insurance: Mathematics and
      Economics}, 44(2), 182–198.
\item Li, D. X. (2000). On default correlation: a copula function approach.
      \textit{Journal of Fixed Income}, 9(4), 43–54.
\item Companion notebook: \texttt{copulas-tail-dependence.ipynb} — reproduces
      every figure from raw data (seed 11).
\end{notereferences}

\end{document}
